\section{All tensor counts for the two ordinary alternating labels}
\label{sec:alternating-two-tail-all-tensors}

\paragraph{Status and exact scope.}
This is a separate new derivation, with all new proofs written below.
It leaves the reviewed D-ATT-01--06 unchanged.  The complete root and /root/imprimitive whole-proof/code
readings have passed, with a distinct matrix supplement and
explicit disclosure of the shared preceding framework.  The original imaginary quadratic field is non-Eisenstein;
it is only assumed to be contained in the character field of the
actual inducing representation.  Equality of those fields is not
assumed.  The result is an obstruction to an actual period-rank-nine
curve source, not a dimension-based construction of a source.

Fix $n\ge15$.  Let $W$ be one of the two ordinary rational
$A_n$-modules obtained by restricting
\[
 S^{(n-2,2)},\qquad S^{(n-2,1,1)}=\Lambda^2\operatorname{Std}_n.
 \tag{ATX.1}
\]
Put $a=\dim W$, and use the constants
\[
\begin{array}{c|c|c|c}
 W&a&d_0&b_0\\ \hline
 S^{(n-2,2)}&n(n-3)/2&n-3&1\\
 \Lambda^2\operatorname{Std}_n&(n-1)(n-2)/2&n-2&\lfloor(n-1)/2\rfloor
\end{array}
\qquad A=\frac{n-5}{2},\quad B=\frac n4.
\tag{ATX.2}
\]
For an operator of degree $v$ write
$\nu(T)=v-\max_z\dim\ker(T-zI)$.  For $x\in S_n$ let
$\operatorname{ind}(x)=n-\#\{\text{cycles of }x\}$, including fixed
points among its cycles.  The complete previously reviewed UP/EX
proofs, collected in D-ATT-01 and D-ATT-04, give
\[
 \nu(Wx)\ge d_0,\qquad
 \nu(Wx)\ge A\operatorname{ind}(x)+B\quad(x\ne1),
 \qquad \operatorname{tr}(Wx)\ge-b_0.
 \tag{ATX.3}
\]
The two support inequalities apply to every target eigenvalue
when $x\ne1$; neither is asserted for the identity.  The trace
bound holds for all $x$.  This distinction is retained throughout.

Let a faithful finite inducing image $J$ act on
$V=W^{\otimes t}$, of degree $e=a^t$, and normalize the ordinary
core $A_n^t$, acting transitively on its $t$ factors.  Here $t\ge3$.
Its compatible normalizer is scalar times the rational tensor-wreath
image of $S_n\wr S_t$.  Indeed $\operatorname{Aut}(A_n)=S_n$ in
this range, and after implementing its factor automorphisms and
factor permutation the remaining operator centralizes the irreducible
core and is scalar.  That rational wreath image meets the scalars
trivially: a scalar centralizes every individual $A_n$ factor, so the
factor permutation is trivial and each $S_n$ automorphism centralizes
$A_n$, hence is the identity.  Thus
\[
                  g=\lambda(g)T(g)
 \tag{ATX.4}
\]
is unique, and both factors in (ATX.4) are actual homomorphisms.
The rational image is faithful as a tensor-wreath action by the
same centralizer argument.  In particular discarding the scalar
gives the actual permutation action on $nt$ letters in $t$ blocks.

\subsection*{D-ATT-07. Two moving tensor-factor branches}

\begin{lemma}[D-ATT-07]
Suppose the actual inducing variation has period rank nine and its
character field contains an imaginary quadratic field
$K\ne\mathbf Q(\sqrt{-3})$.  Then exactly two genuine inertia
entries move the tensor factors.  Their permutations are inverse
full $t$-cycles.
\end{lemma}

\begin{proof}
The reviewed polarized Shapiro/Hurwitz factorization applies to
the actual intermediate connected finite cover, without changing
the inducing field.  Its six-dimensional multiplicity gives
$6=2e(h-1)+D$, where $h$ is the intermediate base genus and
$D$ is the sum of actual moving ranks.  If $h\ge1$, the finite
marked-cover dimension is less than nine: $h=1$ gives at most six
genuine marks, and $h\ge2$ is impossible since $e>3$.
Consequently
\[
       h=0,\qquad r\ge12,\qquad D=2e+6,
 \tag{ATX.5}
\]
for a genuine generating product-one inertia tuple.  These are
necessary conditions from the actual variation, not sufficient
conditions for dominance.

Let $\lambda(J)=\mu_M$.  Rationality of $T(g)$ gives
$\mathbf Q(\chi_J)\subseteq\mathbf Q(\zeta_M)$.  If
$\lambda(J)\subseteq\mu_6$, this field could not contain $K$.
The quotient by $\lambda(J)\cap\mu_6$ is therefore nontrivial.
Generation and product one force at least two inertia entries with
$\lambda\notin\mu_6$.

We recall the complete scalar estimate of D-ATT-02.  The trace of
each power of a rational tensor-wreath operator is a product of
cycle-product traces.  If negative, one factor is at least $-b_0$
and all other absolute values are at most $a$.  Thus, with
$c=b_0/a$, every power has trace at least $-ce$.
For a primitive root of order outside $\{1,2,3,4,6\}$, rationality
makes its multiplicity at most $e/4$ by its cyclotomic conjugates.
For a fourth root, equal multiplicities at $i,-i$ and
$\operatorname{tr}(T^2)\ge-ce$ give multiplicity at most
$(1+c)e/4$.  Therefore each of the two indicated actual scalar
marks has moving rank at least
\[
                         \frac{3-c}{4}e.
 \tag{ATX.6}
\]

Every genuine inertia has moving rank at least $\mu e$, where
$\mu=d_0/a$.  If it fixes the factors and one of its within-factor
permutations is nonidentity, this follows from (ATX.3) and convolution of
eigenspace multiplicities.  A genuine scalar has full moving rank.
If it moves factors, restrict to one nontrivial tensor cycle of
length $l$.  After diagonalizing its cycle product, the constant
words number $a$, and all other word orbits have length at least
two.  Each word orbit contributes at most one vector to a specified
eigenspace.  Its largest eigenspace fraction is at most
$(1+a^{1-l})/2\le(1+1/a)/2$.  Convolution with the remaining
cycles preserves that bound.  Thus every factor-moving inertia
has moving rank at least
\[
               s e,\qquad s=\frac{a-1}{2a}\ge\mu.
 \tag{ATX.7}
\]
The last inequality follows directly from $2d_0\le a-1$ in
(ATX.2), for $n\ge15$.

If at least three inertia entries move factors, select the two
marks in (ATX.6) and a third, distinct factor-moving entry.
Together with at least nine other genuine entries, they give
\[
 D\ge\left(2\frac{3-c}{4}+s+9\mu\right)e
   =2e+a^{t-1}\left(9d_0-\frac{b_0+1}{2}\right)>2e+6.
 \tag{ATX.8}
\]
Here $b_0\le d_0$, $d_0\ge12$ and $a\ge90$.  Thus at most two
entries move factors.  Transitivity on $t>1$ factors requires at
least one, and product one forbids exactly one.  There are exactly
two, their permutations are inverse, and their cyclic subgroup
must be transitive.  Hence both are full $t$-cycles.  This argument
does not assign a deformation or a separate geometric family to
any scalar mark.
\end{proof}

\subsection*{D-ATT-08. A weighted tensor-cycle estimate}

\begin{lemma}[D-ATT-08]
Let $X$ be any invertible diagonalizable degree-$a$ operator, and let $T$ be
a weighted tensor $t$-cycle whose cycle product is $X$.
Let $C_d(a)$ denote the number of primitive cyclic word orbits
of length $d$ on an alphabet of size $a$, and put
\[
 R_t(a)=\sum_{\substack{d\mid t\\d<t}}(t-d)C_d(a).
 \tag{ATX.9}
\]
For every nonzero scalar $\lambda$,
\[
 \operatorname{rank}(1-\lambda T)
 \ge\left(1-\frac1t\right)a^t
       +\frac{a^{t-1}}t\nu(X)-\frac{R_t(a)}t.
 \tag{ATX.10}
\]
For $a\ge4$ and $t\ge3$,
\[
                             R_t(a)\le2a^{t-2}.
 \tag{ATX.11}
\]
\end{lemma}

\begin{proof}
Successive changes of basis in the tensor factors conjugate a
weighted cycle to one with $X$ on one factor and the identity on
the others.  For example, if the cycle sends
$(v_1,\ldots,v_t)$ to $(X_1v_t,X_2v_1,\ldots,X_tv_{t-1})$,
take change-of-basis factors $I,X_2,X_3X_2,\ldots,X_t\cdots X_2$.
The cycle product is $X_1X_t\cdots X_2$, up to the harmless
choice of starting factor.  Its $t$th power applies $X$ to every
factor.

Choose an eigenbasis of $X$, with eigenvalues
$\alpha_1,\ldots,\alpha_a$, repeated with their multiplicities.
For a proposed eigenvalue $z$ of $T$, let $N_z$ count all ordered
words $(i_1,\ldots,i_t)$ satisfying
$\prod_j\alpha_{i_j}=z^t$.  Fixing the first $t-1$ indices shows
\[
                          N_z\le a^{t-1}(a-\nu(X)).
 \tag{ATX.12}
\]
A cyclic word orbit of length $d\mid t$ gives a weighted cyclic
matrix of size $d$.  Its nonzero cyclic weights make each of its
$d$ eigenvalues simple.  If it contributes to the $z$-eigenspace,
all its $d$ words are counted by $N_z$.  Its contribution to
$t\dim\ker(T-zI)-N_z$ is therefore at most $t-d$.
An orbit not contributing gives a nonpositive contribution.
Full-length orbits contribute at most zero.  Summing all shorter
orbits proves
$t\dim\ker(T-zI)\le N_z+R_t(a)$, hence (ATX.10).
Repeated eigenvalues do not merge the basis-word orbits and
cause no change in this counting argument.

For clarity, the same error has the exact fixed-rotation expression
\[
                    R_t(a)=\sum_{j=1}^{t-1}a^{\gcd(t,j)}.
 \tag{ATX.19}
\]
An orbit of length $d<t$ has $d$ words and is fixed by exactly
$t/d-1$ nonidentity rotations, so it contributes $t-d$ to both
counts.  Thus the definition includes the full short-orbit defect,
not merely the number of short words.

For $t=3$, $R_3(a)=2a$.  For $t=4$, the shorter orbits have
length one or two, so $R_4(a)=a^2+2a\le2a^2$ for $a\ge2$.
For $t\ge5$, every proper divisor is at most $\lfloor t/2\rfloor$
and $C_d(a)\le a^d/d$.  Thus
\[
 R_t(a)\le(t-1)\sum_{j=1}^{\lfloor t/2\rfloor}a^j
       <\frac{(t-1)a^{\lfloor t/2\rfloor+1}}{a-1}.
 \tag{ATX.13}
\]
The inequality
$t-1\le2(a-1)a^{\lceil t/2\rceil-3}$ holds at $t=5,6$
for $a\ge4$.  Increasing $t$ by two multiplies the right side
by $a$, while $t+1\le a(t-1)$.  Induction in each parity
therefore proves it for all $t\ge5$.  Substituting it in
(ATX.13) proves (ATX.11).  This is a universal bound; no finite
word enumeration is used to extend its range.
\end{proof}

\subsection*{D-ATT-09. Exclusion for every larger tensor count}

\begin{theorem}[D-ATT-09]
No source of period rank nine in the setting above exists for
any $t\ge3$.  Together with D-ATT-06 and the already reviewed
single-factor UP/EX exclusions, all transitive ordinary core
tensor counts for the two labels (ATX.1), $n\ge15$, are excluded
in the non-Eisenstein original-field problem.
\end{theorem}

\begin{proof}
Use the actual inertia tuple of (ATX.5).  By D-ATT-07 there are
exactly two full tensor cycles.  All other entries fix the
factors.  For such an entry, write its rational part as
$Wx_1\otimes\cdots\otimes Wx_t$, and set
$I=\sum_j\operatorname{ind}(x_j)$.  If precisely $b>0$ of the
$x_j$ are nonidentity, (ATX.3) gives
\[
 \max_j\nu(Wx_j)\ge\frac{AI+bB}{b}
                 \ge\frac{AI}{t}+B.
 \tag{ATX.14}
\]
Convolution of eigenspaces therefore bounds its actual moving
rank, with any scalar, below by
\[
                         a^{t-1}\left(\frac{AI}{t}+B\right).
 \tag{ATX.15}
\]
If all $x_j=1$, the genuine entry is a nonidentity scalar, with
moving rank $e$, and (ATX.15) still holds since $I=0$ and $B<a$.
The full $B$ term in (ATX.15), rather than $B/t$, is essential.

Let $h_1,h_2\in S_n$ be the products around the two full cycles.
For both the identity and nonidentity products, (ATX.3) gives the
valid weaker bound $\nu(Wh_i)\ge A\operatorname{ind}(h_i)$.
D-ATT-08 therefore bounds the sum of the two cycle moving ranks
below by
\[
 2\left(1-\frac1t\right)e
  +\frac{Aa^{t-1}}t\bigl(\operatorname{ind}(h_1)
                         +\operatorname{ind}(h_2)\bigr)
  -\frac{2R_t(a)}t.
 \tag{ATX.16}
\]
No positive $B$ term is demanded at an identity cycle product.
No restriction on whether the actual scalar is a $t$th root
is needed in this estimate.

The scalar-forgetting homomorphism (ATX.4) gives an actual
transitive action on $nt$ letters.  The core is transitive in
each block and the full cycle is transitive on the blocks.
Take the inverse image of a letter stabilizer in the actual
connected $J$-cover of the sphere.  Its quotient is a connected
degree-$nt$ cover.  A factor-fixing entry has permutation index
$I$, whereas a full tensor cycle has permutation index
$n(t-1)+\operatorname{ind}(h_i)$: a cycle of length $l$ of $h_i$
becomes a cycle of length $tl$ on the letters.
Riemann--Hurwitz on this actual quotient consequently gives
\[
 E:=\sum_{\text{factor-fixing entries}} I
       +\operatorname{ind}(h_1)+\operatorname{ind}(h_2)
       \ge2n-2.
 \tag{ATX.17}
\]
Scalar inertia contributes zero to this permutation index but
is retained among the genuine coefficient-cover marks.

There are at least ten factor-fixing entries.  Summing
(ATX.15)--(ATX.16), then using (ATX.11) and (ATX.17), yields
\[
\begin{split}
 D&\ge2e+\frac{a^{t-1}}t
       \left[A(2n-2)+10tB-2a-\frac4a\right]\\
  &=2e+\frac{a^{t-1}}t
       \left[\left(\frac{5t}{2}-3\right)n
                    +\kappa-\frac4a\right],
 \qquad
 \kappa=\begin{cases}5,&W=S^{(n-2,2)},\\3,&W=\Lambda^2\operatorname{Std}_n.
              \end{cases}
\end{split}
\tag{ATX.18}
\]
For $t\ge3$, $5t/2-3\ge3t/2$, and
$\kappa-4/a>0$ since $a\ge90$.  The excess in (ATX.18) is
therefore greater than $3na^{t-1}/2>6$, contradicting (ATX.5).
This proves the theorem for all $n\ge15$ and all $t\ge3$.
\end{proof}

\paragraph{Checks, sources and remaining scope.}
The new word-orbit and all-size source arguments are universal
proofs.  A separate exact supplement checks actual weighted cyclic
matrices, short-orbit errors, permutation quotient indices and
endpoint arithmetic.  It is not a tuple or Hurwitz-component
enumeration.  The scalar-field and ordinary-label inputs are the
complete reviewed D-ATT-01--04 proofs; no additional external
classification or numerical premise is introduced.  Actual
intermediate-cover factorization and Riemann--Hurwitz are retained
explicitly.  Other ordinary alternating labels, spin modules,
nontransitive mixtures and arbitrary nonalternating cores are not
decided here.  No positive supplier or algebraic cycle is claimed.
